\documentclass[11pt]{article}
\usepackage{fullpage}
%% color for the links 
\usepackage[usenames,dvipsnames]{color}
\usepackage{longtable, graphicx}
\usepackage{ebgaramond}		% for boldface small caps
\usepackage{changepage}		% for indented paragraph blocks
% We will generate all images so they have a width \maxwidth. This means
% that they will get their normal width if they fit onto the page, but
% are scaled down if they would overflow the margins.
\makeatletter
\def\maxwidth{\ifdim\Gin@nat@width>\linewidth\linewidth
\else\Gin@nat@width\fi}
\makeatother
\let\Oldincludegraphics\includegraphics
% \renewcommand{\includegraphics}[1]{\Oldincludegraphics[width=\maxwidth]{#1}}
%% hyperlinks
\usepackage[
	colorlinks=true,
	urlcolor=MidnightBlue,
	plainpages=false,
]{hyperref}% color for the links 
\setlength{\parindent}{0pt}
\setlength{\parskip}{6pt plus 2pt minus 1pt}
\setlength{\emergencystretch}{3em}  % prevent overfull lines
\setcounter{secnumdepth}{0}

\begin{document}

\begin{minipage}[c]{0.4\linewidth}
  \includegraphics[height = 2cm]{../images/rice_smgt.png}
\end{minipage}
\begin{minipage}[c]{0.6\linewidth}
  \raggedleft
  {\bf Course:} SMGT $373/374$ -- Sport Analytics Internship
\end{minipage}

\hrule

~\\
\textbf{\textsc{Instructor Contact Information}}
\begin{adjustwidth}{1cm}{0cm}
  \begin{tabular}{lll}
  {\bf Instructor:} Scott Powers                                                      & \hspace{1cm}  & {\bf Administrator:} Kerri Barber\\
  {\bf Office:} Kraft $246$                                                           &               & {\bf Office:} Kraft $249$\\
  {\bf Email:} \href{mailto:scott.powers@rice.edu}{scott.powers@rice.edu}\\
  {\bf Office Hours:} \href{https://saberpowers.com/bookme}{saberpowers.com/bookme}   &               & {\bf Office Hours:} by appointment
  \end{tabular}
\end{adjustwidth}

~\\
\textbf{\textsc{Required Texts and Materials}}
\begin{adjustwidth}{1cm}{0cm}
	{\bf Class Website:} canvas.rice.edu\\
	{\bf Required Text:} none\\
	{\bf Required Material:} none
\end{adjustwidth}

~\\
\textbf{\textsc{Requirements}}
\begin{adjustwidth}{1cm}{0cm}
  \begin{enumerate}
    \item You must secure an internship in sport analytics and identify your internship supervisor before registering for the course. The instructor must approve the caliber of the internship.
    \item You must complete at least $150$ working hours before the last day of classes for the semester.
  \end{enumerate}
\end{adjustwidth}

~\\
\textbf{\textsc{Registration}}
\begin{adjustwidth}{1cm}{0cm}
  \begin{enumerate}
  \item For permission to register, send the instructor an email satisfying the following criteria:
  \begin{itemize}
    \item You have CC'd your internship supervisor on the email.
    \item You have attached a Job Description (JD) as a PDF file with a list of job responsibilities with percentage weights that sum to $100\%$. If your JD does not include weights, ask your internship supervisor to help you add them to the JD.
    \item You have attached a Supervisor Acknowledgment Form completed by your internship supervisor. \href{https://saberpowers.com/teaching/smgt373/acknowledgment.pdf}{saberpowers.com/teaching/smgt373/acknowledgment.pdf}
  \end{itemize}

  \item After sending the email, submit a Course Registration Override request following the instructions at \href{https://registrar.rice.edu/course-registration-override-request-process}{registrar.rice.edu/course-registration-override-request-process}.

  \item After the instructor approves your Course Registration Override, you must still log in to ESTHER and officially add the class to your schedule.

  \item After registering, you can update your variable credit hours (if necessary) following the instructions at \href{https://registrar.rice.edu/students/updating_variable_credit}{registrar.rice.edu/students/updating\_variable\_credit}.
  \end{enumerate}
\end{adjustwidth}

~\\
\textbf{\textsc{Course Objectives and Learning Outcomes}}
\begin{adjustwidth}{1cm}{0cm}
  After successfully completing this course, you will have accomplished the following:
  \begin{itemize}
    \item Performed an internship at a site where you obtained experience to help advance your career in sports.
    \item Completed the internship with clearly defined goals and expectations, allowing you to list specific accomplishments on your resume that can be cited to demonstrate competency to future employers.
  \end{itemize}
\end{adjustwidth}

\newpage

\textbf{\textsc{Expectations}} \hspace{1cm} {\it Credit: Tom Stallings}
\begin{adjustwidth}{1cm}{0cm}

  Do not assume you can take off work if you have a friend in town, a class project you are behind on or your birthday, etc. These are not acceptable reasons to take the day off in an internship. In addition, internship site obligations may supersede university declarations. For example, Rice may announce school holidays when classes are off, but unless requested in advance, internship obligations would be expected to be maintained.

  Please arrange in advance any days that you will be off. Family reunions, weddings, varsity athletic obligations, club sports obligations or other Rice-sponsored organization activities (e.g. Rice Speech and Debate Club competitions, theater or music performances, conferences) are examples of acceptable reasons to request time off. In a business setting, you will have other obligations that may take you away from the workplace. The key is to let your supervisor know as far in advance as possible. You and your supervisor must agree in advance to these requested days off, and it is recommended that you submit these dates in writing to your site supervisor with a copy sent to your instructor. If you have any doubt, reach out to your instructor in advance of making the request with your employer.

  Dress appropriately! Part of showing respect for your internship is showing respect for the dress code in your office. Many offices frown upon jeans, sandals, t-shirts, revealing clothing and other attire common in the college setting. If you are not sure what appropriate office attire is for your internship, take a look around at your colleagues or consult directly with your site supervisor or your company's HR department.

  To really make an impact, an intern should seek to identify an organizational need, research ways to address that need, propose solutions to solve the challenges, review the proposal with your site supervisor and implement the recommendations successfully. Your goal is to make a substantive impact, which will increase your chances of securing a strong recommendation for future opportunities or receiving a full-time job offer.
\end{adjustwidth}

\textbf{\textsc{How Your Grade Is Determined}}
\begin{adjustwidth}{1cm}{0cm}
  Your grade is based on your percentage earned of available points.
  \begin{center}
    \begin{tabular}{rclcrclcrcl}
                      &               &     & & $[93\%,\,\infty)$ & $\rightarrow$ & A & & $[90\%, 93\%)$ & $\rightarrow$	& A--\\
      $[87\%, 90\%)$  & $\rightarrow$	& B+  & & $[83\%, 87\%)$  & $\rightarrow$ & B & & $[80\%, 83\%)$ & $\rightarrow$	& B--\\
      $[77\%, 80\%)$  & $\rightarrow$	& C+  & & $[73\%, 77\%)$  & $\rightarrow$ & C & & $[70\%, 73\%)$ & $\rightarrow$	& C--\\
      $[67\%, 70\%)$  & $\rightarrow$	& D+  & & $[63\%, 67\%)$  & $\rightarrow$ & D & & $[60\%, 63\%)$ & $\rightarrow$  & D--\\
                      &               &     & & $[ 0\%, 60\%)$  & $\rightarrow$ & F\\
    \end{tabular}
  \end{center}
  ~\\
  \textbf{Hour Logs.} ($50$pts) There are $5$ hour logs, one every three weeks. Your score on each hour log is equal to [the number of hours you log] divided by [the number of credit hours you enrolled]. For example, if you enrolled in $3$ credit hours and log $30$ hours, that is worth $10$pts. Hour log points are capped at $50$ total.

  \textbf{Experience Picture \& Caption.} ($5$pts) This is for marketing purposes (graded for completion).

  \textbf{Midterm Performance Evaluation.} ($25$pts) It is your responsibility to remind your internship supervisor to complete your midterm performance evaluation halfway through the semester. Your score will be the weighted average of your performance on each responsibility in the JD (weighted according to JD weights), which will Need Improvement ($70\%$), Meet Expectations ($85\%$) or Exceed Expectations ($100\%$).

  \textbf{Final Performance Evaluation.} ($25$pts) It is your responsibility to remind your internship supervisor to complete your final performance evaluation at the end of the semester. Your score will be the weighted average of your performance on each responsibility in the JD (weighted according to JD weights), which will Need Improvement ($70\%$), Meet Expectations ($85\%$) or Exceed Expectations ($100\%$).

  \textbf{Late Policy.} Late submissions are subject to a $1$pt deduction per day.
\end{adjustwidth}

\newpage
\textbf{\textsc{FAQ}}
\begin{adjustwidth}{1cm}{0cm}
  \textit{Can I receive credit for an internship I already completed?}\\
  Students may not receive academic credit for a retroactive internship because we have no mechanism to verify the rigor or academic relevance of experiences that were not previously reviewed. To earn academic credit, students must receive active supervision from both the instructor and an approved internship supervisor throughout their internship. Academic credit is not awarded simply for working a set number of hours.\\

  \textit{What qualifies as an internship in sport analytics?}\\
  At least $50\%$ of the job responsbility weight needs to be on sport analytics, which we define as the quantitative analysis of in-game sports data for measuring, predicting, and improving team and individual performance in competitions. Data engineering, sport science, and sport business analytics job functions are acceptable. Coaching, scouting, and finance job functions do satisfy this requirement because they do not involve sufficient quantitative analysis.
\end{adjustwidth}


\newpage
\textbf{\textsc{Rice Honor Code}}
\begin{adjustwidth}{1cm}{0cm}
  In this course, all students will be held to the standards of the Rice Honor Code, a code that you pledged to honor when you matriculated at this institution. If you are unfamiliar with the details of this code and how it is administered, you should consult the Honor System Handbook at \url{http://honor.rice.edu/honor-system-handbook/}. This handbook outlines the University's expectations for the integrity of your academic work, the procedures for resolving alleged violations of those expectations, and the rights and responsibilities of students and faculty members throughout the process.
\end{adjustwidth}

~\\
\textbf{\textsc{Disability Resource Center}}
\begin{adjustwidth}{1cm}{0cm}
  The Americans with Disabilities Act ensures that all qualified individuals have equal access to educational opportunities, regardless of disability. Students with a documented disability who require academic accommodations should contact the Student Disability Resource Center (\href{https://sdrc.rice.edu}{sdrc.rice.edu}) to submit appropriate documentation or confirm that their records are current. The SDRC is located in Allen Center 111 and can be reached at \href{mailto:sdrc@rice.edu}{sdrc@rice.edu} or x$5841$.

  For additional information regarding access and accommodations for individuals with disabilities, please review Rice University Policy 402: \href{https://policy.rice.edu/402}{policy.rice.edu/402}.
\end{adjustwidth}

~\\
\textbf{\textsc{Mental Health Policy}}
\begin{adjustwidth}{1cm}{0cm}
	The wellbeing and mental health of students is important; if you are having trouble completing your coursework, please reach out to the Wellbeing and Counseling Center. Rice University provides cost-free mental health services through the Wellbeing and Counseling Center to help you manage personal challenges that threaten your personal or academic well-being. If you believe you are experiencing unusual amounts of stress, sadness, or anxiety, the Student Wellbeing Office or the Rice Counseling Center may be able to assist you. The Wellbeing and Counseling Center is located in the Gibbs Wellness Center and can be reached at 713-348-3311 (available 24/7).
\end{adjustwidth}

~\\
\textbf{\textsc{Title IX Responsible Employee Notification}}
\begin{adjustwidth}{1cm}{0cm}
  At Rice University, unlawful discrimination in any form, including sexual misconduct, is prohibited under Rice Policy on Harassment and Sexual Harassment (Policy 830) and the Student Code of Conduct.

  Please be aware that all employees of Rice University are ``mandatory reporters'', which means that if you tell me about a situation involving sexual harassment, sexual assault, dating violence, domestic violence, or stalking, I must share that information with the Title IX Coordinator.

  Although I have to make that notification, you will control how your case will be handled, including whether or not you wish to pursue a formal complaint. Our goal is to make sure you are aware of the range of options available to you and have access to the resources you need.

  To report sexual harassment, please contact the Title IX Coordinator at titleix@rice.edu. To explore supportive measures and other resources that are available to you, please visit the Office if Interpersonal Misconduct Prevention and Support at safe.rice.edu.
\end{adjustwidth}

~\\
\textbf{This syllabus is only a guide for the course and is subject to change with advance notice.}

\end{document}
